% Part: incompleteness
% Chapter: arithmetization-syntax
% Section: introduction

\documentclass[../../../include/open-logic-section]{subfiles}

\begin{document}

\olfileid{inc}{art}{int}
\olsection{تعارف}

قابلِ حسابیت تے منطق نوں جوڑن لئی سانوں ایسا طریقہ چاہیدا اے
جیہدے نال منطق دیاں چیزاں (علامتاں، ٹرم، !!{formula}s تے
!!{derivation}s)، اوہناں اُتے ہوندے عمل، تے اوہناں دیاں خاصیتاں
تے تعلقاں نوں حسابی طریقیاں نال ورت سکّیے۔ ایہ کم سیدھا وی ہو سکدا
اے، جے اسیں علامتاں، علامتاں دے سلسلیاں تے اوہناں توں بنیاں ہور
چیزاں اُتے قابلِ حساب فنکشن تے تعلق ویکھیے۔ منطقی نحو دی ہر شے
متناہی اے تے علامتاں دے !!a{enumerable} مجموعیاں توں بندی اے،
ایس لئی حساب دے کجھ نمونیاں وچ ایہ طریقہ چل سکدا اے۔ پر حساب دے
ہور نمونے، جیویں بازگشتی فنکشن، صرف عدداں، عدداں وچلے تعلقاں
تے عدداں اُتے فنکشناں تک محدود نیں۔ آخرکار اسیں کجھ نظرییاں دے
اندر وی نحو نال کم کرنا چاہندے آں، خاص کر اوہناں نظرییاں وچ
جیہڑے حساب دی زبان وچ لکھے گئے نیں۔ اوتھے نحو نوں \emph{عددی
روپ دینا} ضروری اے۔ یعنی نحوی چیزاں نوں عدداں دے روپ وچ، اوہناں
اُتے ہوندے عملاں نوں حسابی فنکشناں دے روپ وچ تے اوہناں دے تعلقاں
نوں حسابی تعلقاں دے روپ وچ پیش کرنا پیندا اے۔ نحوی چیزاں نوں
عدد دین دا ایہ خیال لائبنز دے ویلے توں چلیا آؤندا اے۔

علامتاں نوں «کوڈ» وجوں عدد دینا نسبتاً سوکھا اے۔ کجھ علامتاں
لئی ذرا اوکھ اے، کیوں جے، مثلاً، !!{variable}s لامتناہی نیں تے
دلیلاں دی ہر مقررہ تعداد~$n$ والے !!{function}s وی لامتناہی نیں۔
فیر وی علامتاں نوں ترتیب نال عدد دِتے جا سکدے نیں، جیویں
$\Obj v_2$ تے $\Obj v_3$ نوں وکھو وکھ کوڈ۔ علامتاں دے سلسلیاں
(جیویں ٹرماں تے !!{formula}s) نوں کوڈ دینا ہور اوکھا اے۔ پر جے
عدداں دے سلسلیاں نال خالص حسابی طریقے نال کم ہو سکدا اے
(مثلاً اولی عدداں دیاں قوتاں نال سلسلیاں نوں کوڈ دے کے)، تاں
اسیں اک اک علامت دی کوڈکاری توں اگے ودھ کے علامتاں دے سلسلیاں
نوں وی کوڈ دے سکدے آں۔ ایس توں اگے !!{formula}s دے سلسلیاں یا
ہور طرح دیاں بناوٹاں، جیویں !!{derivation}s، نوں وی کوڈ دے سکدے
آں۔ ایس وادھو کوڈکاری نوں «گوڈل نمبر دینا» (Gödel numbering)
آکھدے آں۔ ہر ٹرم، !!{formula} تے !!{derivation} نوں اک گوڈل
نمبر دِتا جاندا اے۔

جے اسیں علامتاں دے سلسلیاں دی تھاں اوہناں دیاں علامتاں دے
کوڈاں دے سلسلے لئیے، تے سلسلیاں نوں کوڈ دین دا ایسا طریقہ چنیے
جیہڑے نال قابلِ حساب فنکشناں راہیں کم ہو سکے، تاں گوڈل نمبراں
نال وی قابلِ حساب فنکشناں راہیں کم ہو سکدا اے۔ عمل وچ سارے
متعلقہ فنکشن بنیادی بازگشتی ہون گے۔ مثلاً سلسلے دے کوڈ توں
اوہدی لمبائی لبھن والا فنکشن تے اوہدا $i$-واں رکن کڈھن والا
فنکشن، دونوں بنیادی بازگشتی نیں۔ جے سلسلے دا کوڈ، مثلاً،
!!a{formula}~$!A$ دا گوڈل نمبر ہووے، تاں ایس توں فوراً پتا لگدا
اے کہ !!a{formula} دی لمبائی تے !!a{formula} دے $i$-ویں نشان
دا کوڈ وی $!A$ دے گوڈل نمبر توں کڈھیا جا سکدا اے۔ ایہ ثابت
کرنا ذرا اوکھا اے کہ، مثلاً، صحیح طرح بنے ٹرم یا درست
!!{derivation} دا گوڈل نمبر ہون دی خاصیت بنیادی بازگشتی اے۔
فیر وی ایہ ثابت ہو سکدا اے، کیوں جے متعلقہ سلسلیاں (ٹرم،
!!{formula}s تے !!{derivation}s) دیاں تعریفاں استقرائی نیں۔

مثال لئی بدل دا عمل ویکھو۔ جے $!A$ اک فارمولا، $x$ اک متغیر
تے $t$ اک ٹرم ہووے، تاں $\Subst{!A}{t}{x}$ اوہ نتیجہ اے جیہڑا
$x$ دیاں~$!A$ وچ ساریاں آزاد آمدواں دی تھاں $t$ رکھ کے ملدا اے۔
ہن فرض کرو کہ $!A$، $x$ تے $t$ نوں ایس ترتیب نال گوڈل نمبر
$k$، $l$ تے $m$ دِتے گئے نیں۔ ایہو کوڈکاری $\Subst{!A}{t}{x}$
نوں وی اک گوڈل نمبر دیندی اے، من لو~$n$۔ $k$، $l$ تے $m$ نوں
$n$ ول لیجان والا ایہ فنکشن بدل دے عمل دا حسابی روپ اے۔ جدوں
بدل دا عمل $!A$، $x$ تے $t$ توں $\Subst{!A}{t}{x}$ دیندا اے،
تاں بدل دا عددی روپ گوڈل نمبراں $k$، $l$ تے $m$ توں گوڈل
نمبر~$n$ دیندا اے۔ اگے اسیں ویکھاں گے کہ ایہ فنکشن بنیادی
بازگشتی اے۔

نحو نوں عددی روپ دینا صرف نظری دلچسپی دی گل نہیں۔ البتہ شروع
وچ ایہ سمجھنا سوکھا نہیں سی کہ حساب دی زبان ورگیاں زباناں وی
عبارتاں دیاں پیچیدہ خاصیتاں نوں رسمی روپ دے سکدیاں نیں، حالانکہ
اوہناں وچ زباناں بارے «گل کرن» دا کوئی اپنا وسیلہ نہیں۔ ایس توں
اگلا سوال ایہ اے کہ ایس زبان وچ لکھیا کوئی نظریہ، جیویں پیانو
حساب، اپنی زبان بارے کیہ \emph{ثابت} کر سکدا اے (مثلاً، کیہ
!!{sentence}s ثابت کیتے جا سکدے نیں، یا کیہ اوہ سچ نیں)۔ ایس
سوال توں گوڈل تے تارسکی دے مشہور قضیاں تک پہنچدے آں، جیہڑے
ثبوت تے تعریف دیاں حداں دسدے نیں۔ گوڈل دے قضیے ثابت نہ ہو
سکّن بارے نیں تے تارسکی دا قضیہ ایس گل بارے اے کہ سچ دی تعریف
نہیں دِتّی جا سکدی۔ نحو نوں عددی روپ دین دا طریقہ قابلِ حسابیت
دے نظریے دے کجھ اہم نتیجے ثابت کرن لئی وی کم آؤندا اے، جیویں
نظرییاں دی حسابی طاقت یا حساب دے وکھو وکھ نمونیاں دے آپسی تعلق
بارے نتیجے۔ ایہ طریقہ ہور چیزاں تے خاصیتاں نوں عددی روپ دین دا
نمونہ وی اے۔ مثلاً ٹیورنگ مشیناں دیاں مکمل حالتاں تے اوہناں
دیاں حسابی کارروائیاں نوں وی ایویں عددی روپ دِتا جا سکدا اے۔
ایس نال اک نمونے (مثلاً ٹیورنگ مشیناں) دیاں حسابی کارروائیاں
دوجے نمونے (مثلاً بازگشتی فنکشناں) وچ نقل کر کے چلائیاں جا
سکدیاں نیں۔

\end{document}
